%%%%%%%%%%%%%%%%%%%%%%%%%%%%%%%%%%%%%%%%%%%%%%%%%%%%%%%
% NOTAS TEORIA DE TRANSPORTE OPTIMO
%
% Contains ?? figures. 
% File started: ??
% First circulated version: ??
% Last version:
% Last modification: -
%%%%%%%%%%%%%%%%%%%%%%%%%%%%%%%%%%%%%%%%%%%%%%%%%%%%%%%

\documentclass[11pt,a4paper]{amsart}
\usepackage{xr-hyper}
\newcounter{chapter}

\usepackage{amssymb}
\usepackage{hyperref}
\usepackage{listings}
\usepackage{mathtools}
\usepackage{enumitem}

% Fuzz -------------------------------------------------------------------
\hfuzz2pt % Don't bother to report over-full boxes if over-edge is < 2pt
% Line spacing -----------------------------------------------------------

%%%%%%%%%fin pegada%%%%%%%%%%%%%%%%%%5

\usepackage[T1]{fontenc}
\usepackage[spanish]{babel}
\usepackage{graphicx}
\usepackage{tikz}
\usetikzlibrary{babel}

\setlength{\textwidth}{15.5cm} 
\setlength{\textheight}{24cm}
\setlength{\oddsidemargin}{0cm}
\setlength{\evensidemargin}{0cm}


\parskip 3pt

%%
%% OPERADORES
%%
\renewcommand{\Re}{\operatorname{Re}}
\newcommand{\tr}{\operatorname{tr}}
\newcommand{\id}{\operatorname{id}}
\def\div{\operatorname {\mathrm{div}}}
\def\Dom{\operatorname {\mathrm{Dom}}}
\def\argmin{\operatorname {\mathrm{arg\, min}}}
\def\argmax{\operatorname {\mathrm{arg\, max}}}
\def\gen{\operatorname {\mathrm{gen}}}
\def\Im{\operatorname {\mathrm{Im}}}
\def\Nu{\operatorname {\mathrm{Nu}}}
\def\sop{\operatorname {\mathrm{sop}}}
\def\diam{\operatorname {\mathrm{diam}}}
\def\dist{\operatorname {\mathrm{dist}}}
\def\var{\operatorname {\mathrm{Var}}}
\def\codim{\operatorname {\mathrm{codim}}}
\def\pint{\operatorname {--\!\!\!\!\!\int\!\!\!\!\!--}}

%%
%% EPSILON
%%
\newcommand{\ve}{\varepsilon}

%%
%% LETRAS MATHBB
%%
\newcommand{\R}{{\mathbb R}}
\newcommand{\Z}{{\mathbb Z}}
\newcommand{\C}{{\mathbb C}}
\newcommand{\Q}{{\mathbb Q}}
\newcommand{\N}{{\mathbb N}}
\newcommand{\E}{{\mathbb E}}
\newcommand{\V}{{\mathbb V}}
\renewcommand{\P}{{\mathbb P}}
\newcommand{\II}{{\mathbb I}}
\newcommand{\OO}{{\mathbb O}}
\newcommand{\Sb}{{\mathbb S}}
\newcommand{\Tb}{{\mathbb T}}

%%
%% LETRAS CALIFGRAFICAS
%%
\newcommand{\F}{{\mathcal F}}
\newcommand{\Hc}{{\mathcal H}}
\newcommand{\Tc}{{\mathcal T}}
\newcommand{\Ec}{{\mathcal E}}
\newcommand{\Vc}{{\mathcal V}}
\newcommand{\A}{{\mathcal A}}
\newcommand{\J}{{\mathcal J}}
\newcommand{\G}{{\mathcal G}}
\newcommand{\U}{{\mathcal U}}
\newcommand{\B}{{\mathcal B}}
\newcommand{\Sw}{{\mathcal S}}
\newcommand{\D}{{\mathcal D}}
\newcommand{\M}{{\mathcal M}}
\newcommand{\LL}{{\mathcal L}}

%%
%% LETRAS MATHBF
%%
\newcommand{\ab}{{\mathbf a}}
\newcommand{\bb}{{\mathbf b}}
\newcommand{\cc}{{\mathbf c}}
\newcommand{\ee}{{\mathbf e}}
\newcommand{\ff}{{\mathbf f}}
\newcommand{\gb}{{\mathbf g}}
\newcommand{\hh}{{\mathbf h}}
\newcommand{\n}{{\mathbf n}}
\newcommand{\mm}{{\mathbf m}}
\newcommand{\pp}{{\mathbf p}}
\newcommand{\qq}{{\mathbf q}}
\newcommand{\sbf}{{\mathbf s}}
\newcommand{\tb}{{\mathbf t}}
\newcommand{\uu}{{\mathbf u}}
\newcommand{\vv}{{\mathbf v}}
\newcommand{\ww}{{\mathbf w}}
\newcommand{\xx}{{\mathbf x}}
\newcommand{\yy}{{\mathbf y}}
\newcommand{\zz}{{\mathbf z}}
\newcommand{\I}{{\mathbf 1}}
\newcommand{\Pb}{{\mathbf P}}
\newcommand{\XX}{{\mathbf X}}
\newcommand{\YY}{{\mathbf Y}}

%%
%% LETRAS MATHFRAK
%%

\newcommand{\Gk}{{\mathfrak G}}


%%
%% TEOREMAS
%%
\newtheorem{teo}{Teorema}[section]
\newtheorem{lema}[teo]{Lema}
\newtheorem{prop}[teo]{Proposición}
\newtheorem{cor}[teo]{Corolario}

%%
%% REMARK
%%
\theoremstyle{remark}
\newtheorem{obs}[teo]{Observación}

%%
%% DEFINICION
%%
\theoremstyle{definition}
\newtheorem{defi}[teo]{Definición}
\newtheorem{ejem}[teo]{Ejemplo}
\newtheorem{ejer}[teo]{Ejercicio}
\newtheorem{nota}[teo]{Notación}
\newtheorem{ejercicio}{Ejercicio}[chapter]


%%
%% NUMERACION
%%
\numberwithin{subsection}{section}
\numberwithin{equation}{section}

%%
%% TITULO
%%


\externaldocument[N-]{../notas/Notas-TO}
\newcommand{\cuadernolink}[2]{\noindent\textbf{Cuaderno:} \emph{#1} (\href{https://colab.research.google.com/github/jfbonder/transporte-optimo/blob/main/cuadernos/#2}{abrir en Colab}).\par\medskip}
\pagestyle{plain}

\begin{document}
\renewcommand{\thechapter}{5}
\begin{center}
{\large\sc Teoría de Transporte Óptimo}\\[2pt]
{\small 2do cuatrimestre 2026 --- Julián Fernández Bonder}\\[10pt]
{\LARGE\bf Práctica 5}\\[4pt]
{\large Potenciales de Kantorovich en el caso discreto}
\end{center}
\bigskip
\noindent{\small Los ejercicios son los de la sección de ejercicios del Capítulo~5 de las notas del curso, con la misma numeración. Las referencias a teoremas, proposiciones y secciones remiten a las notas (\url{https://github.com/jfbonder/transporte-optimo}).}
\medskip


\begin{ejercicio}[Un problema $3\times3$ con potenciales no únicos]\label{ejer.discreto.3x3}
Sean $x_i=y_i=i-1$ para $i=1,2,3$, $a=(\tfrac12,\tfrac14,\tfrac14)$,
$b=(\tfrac14,\tfrac14,\tfrac12)$ y $c_{ij}=|x_i-y_j|^2$.
\begin{enumerate}[label=(\alph*)]
\item Hallar un plan óptimo $\pi$ y el costo óptimo. (Conviene ``llenar''
la matriz de arriba a la izquierda hacia abajo a la derecha; en el
Capítulo~\ref{N-cap.1D} se verá por qué esto es óptimo.)
\item Escribir las condiciones de saturación $\varphi_i+\psi_j=c_{ij}$ en el
soporte de $\pi$ y las de admisibilidad fuera de él. Normalizando
$\varphi_1=0$, probar que los pares óptimos forman el segmento
\[
\varphi=(0,\,1-t,\,-t),\qquad\psi=(0,\,1,\,t),\qquad t\in[2,4].
\]
\item Verificar que $\sum_ia_i\varphi_i+\sum_jb_j\psi_j$ no depende de $t$ y
coincide con el costo óptimo.
\item Explicar la no unicidad en términos del grafo bipartito
$G_\pi=\{(i,j):\pi_{ij}>0\}$: ¿cuántas componentes conexas tiene?
\end{enumerate}
\end{ejercicio}

\begin{ejercicio}[Unicidad de los potenciales y conexión del soporte]\label{ejer.discreto.unicidad}
Sea $\pi$ un plan óptimo del problema discreto y $G_\pi$ su grafo bipartito
de soporte.
\begin{enumerate}[label=(\alph*)]
\item Probar que todo par óptimo $(\varphi,\psi)$ satura sobre $G_\pi$
(Ejercicio~\ref{N-ejer.PL.holgura}).
\item Probar que si $G_\pi$ es conexo, los pares óptimos son únicos salvo
$(\varphi+s,\psi-s)$, $s\in\R$.
\item Probar que si $G_\pi$ tiene $k$ componentes conexas, el conjunto de
pares óptimos tiene dimensión a lo sumo $k$, y dar un ejemplo en el que sea
exactamente $k$ (el Ejercicio~\ref{ejer.discreto.3x3} es el caso $k=2$).
\end{enumerate}
\end{ejercicio}

\begin{ejercicio}[Costo $0$--$1$ y variación total]\label{ejer.discreto.01}
Sean $x_1,\dots,x_N$ puntos distintos, $\mu=\sum a_i\delta_{x_i}$,
$\nu=\sum b_i\delta_{x_i}$ (mismo soporte, pesos distintos) y
$c(x,y)=\I_{\{x\ne y\}}$.
\begin{enumerate}[label=(\alph*)]
\item Probar que el costo óptimo es
$\sum_i(a_i-b_i)^+=\tfrac12\sum_i|a_i-b_i|=\|\mu-\nu\|_{TV}$, exhibiendo un
plan que lo alcanza (dejar quieta la masa $\min\{a_i,b_i\}$ en cada punto).
\item Hallar un par óptimo $(\varphi,\psi)$ del problema dual con
$\varphi_i\in\{0,1\}$ y $\psi_i\in\{-1,0\}$, y verificar la dualidad fuerte.
\item Generalizar: para $\mu,\nu\in\mathcal P(X)$ cualesquiera y
$c=\I_{\{x\ne y\}}$, probar que
\[
\min_{\pi\in\Pi(\mu,\nu)}\int c\,d\pi=\|\mu-\nu\|_{TV}.
\]
\end{enumerate}
\end{ejercicio}

\begin{ejercicio}[La cara óptima]\label{ejer.discreto.cara}
Sea $\Pi^0\subset\Pi(a,b)$ el conjunto de planes óptimos del problema
discreto.
\begin{enumerate}[label=(\alph*)]
\item Probar que $\Pi^0$ es un politopo convexo no vacío y que sus puntos
extremales son puntos extremales de $\Pi(a,b)$ (una \emph{cara} del
politopo).
\item Sean $a=b=\frac14(1,1,1,1)$, $x_i=i-1$, $y_j=j$ ($i,j=1,\dots,4$) y
$c_{ij}=|x_i-y_j|$. Probar que el costo óptimo es $1$, exhibir al menos tres
planes óptimos extremales distintos y describir $\Pi^0$.
\item Probar que en (b) todo par óptimo $(\varphi,\psi)$ satura sobre la
unión de los soportes de todos los planes óptimos, y hallar todos los pares
óptimos.
\end{enumerate}
\end{ejercicio}

\end{document}
